\begin{figure}[t]
\resizebox{\linewidth}{!}{%
\begin{tikzpicture}[
    % 定义方框的样式
    box/.style={
        rectangle,
        draw=black,               % 边框颜色
        thick,                    % 边框粗细
        rounded corners=8pt,      % 圆角大小
        fill=gray!15,             % 填充颜色 (浅灰)
        align=center,             % 文字居中
        minimum width=4.2cm,      % 最小宽度
        minimum height=2cm,       % 最小高度
        font=\small,              % 字体大小
        inner sep=5pt
    },
    % 定义连接线的样式
    connector/.style={
        draw=gray!90,             % 线条颜色 (深灰)
        line width=3.5pt          % 线条粗细
    }
]

    % --- 左侧组 (Left Cluster) ---
    
    % 1. 定义中心点坐标 (看不见的锚点，用于定位十字交叉点)
    \coordinate (L_center) at (0, 0);

    % 2. 放置节点
    % Hero/heroine (#1) - 在中心上方
    \node[box] (hero) at ($(L_center)+(0, 2.2)$) {
        \textbf{Hero/heroine}\\
        \#1 (\textit{superior or dominant}\\
        \textit{function})
    };

    % Anima/animus (#4) - 在中心下方
    \node[box] (anima) at ($(L_center)+(0, -2.2)$) {
        \textbf{Anima/animus}\\[1em] % 增加一点行间距
        \#4 (\textit{inferior function})
    };

    % Father/mother (#2) - 在中心左侧
    \node[box] (father) at ($(L_center)+(-3.0, 0)$) {
        \textbf{Father/mother}\\[1em]
        \#2 (\textit{auxiliary function})
    };

    % Puer/puella (#3) - 在中心右侧
    \node[box] (puer) at ($(L_center)+(3.0, 0)$) {
        \textbf{Puer/puella}\\[1em]
        \#3 (\textit{tertiary function})
    };


    % --- 右侧组 (Right Cluster) ---
    
    % 定义右侧组的中心点 (向右偏移 11cm)
    \coordinate (R_center) at (11, 0);

    % Opposing personality (#5)
    \node[box] (opposing) at ($(R_center)+(0, 2.2)$) {
        \textbf{Opposing personality}\\
        \#5 (\textit{same function as \#1}\\
        \textit{but with opposite attitude})
    };

    % Demonic personality (#8)
    \node[box] (demonic) at ($(R_center)+(0, -2.2)$) {
        \textbf{Demonic personality}\\
        \#8 (\textit{same function as \#4}\\
        \textit{but with opposite attitude})
    };

    % Senex/witch (#6)
    \node[box] (senex) at ($(R_center)+(-3.0, 0)$) {
        \textbf{Senex/witch}\\
        \#6 (\textit{same function as \#2}\\
        \textit{but with opposite attitude})
    };

    % Trickster (#7)
    \node[box] (trickster) at ($(R_center)+(3.0, 0)$) {
        \textbf{Trickster}\\
        \#7 (\textit{same function as \#3}\\
        \textit{but with opposite attitude})
    };


    % --- 绘制连接线 (放在背景层，以免遮挡文字框) ---
    \begin{scope}[on background layer]
        % 左侧十字
        \draw[connector] (hero.south) -- (anima.north);
        \draw[connector] (father.east) -- (puer.west);
        
        % 右侧十字
        \draw[connector] (opposing.south) -- (demonic.north);
        \draw[connector] (senex.east) -- (trickster.west);

        % 中间连接桥 (连接两个十字中心)
        % 从左侧十字中心 画到 右侧十字中心
        % 为了美观，我们直接画一条横贯线连接两个内部节点
        %\draw[connector] (puer.east) -- (senex.west);
    \end{scope}

\end{tikzpicture}}

\caption{Archetypal complexes carrying the eight functions of consciousness. \verified}
\label{fig:beebe-fig71}

\end{figure}
